% from ega0-22-fr.tex
\oldpage[0\textsubscript{IV-1}]{198}

Il résulte aussitôt des définitions que le diagramme
\[
  \xymatrix{
    \Theta_{B_{(K)}/K} \ar[r]^{\pi'_{B_{(K)}/K}} \ar[dr]_{\pi_{B_{(K)}/K}}
      & \Omega^1_A\otimes_AK \ar[d]^{\varphi_{K/A}} \\
    & \Omega^1_K
  }
\]
est commutatif, $\pi_{B_{(K)}/K}$ étant l'homomorphisme canonique
correspondant à $i_0$ et $j_0$ (21.3.2.2). Par \emph{restriction} à
$\Xi_{B_{(K)}/K}=\operatorname{Ker}(\pi_{B_{(K)}/K})$, l'homomorphisme
$\pi'_{B_{(K)}/K}$ définit donc un \emph{homomorphisme canonique}
\begin{equation}
\label{0.22.4.5.3-fr}
  \Xi_{B_{(K)}/K}\longrightarrow
  \operatorname{Ker}(\Omega^1_A\otimes_AK\longrightarrow\Omega^1_K)=\Upsilon_{K/A}.
  \tag{22.4.5.3}
\end{equation}
\end{env}

\begin{env}[22.4.6]
\label{0.22.4.6-fr}
Les hypothèses et notations étant celles de (22.4.5), supposons en outre que
$\mathfrak m$ soit un idéal \emph{maximal de carré nul}, donc $K$ un
\emph{corps}, et notons $V$ l'idéal $\mathfrak m$ considéré comme $K$-espace
vectoriel. Comme $K$ est une algèbre formellement lisse sur son corps premier
$\mathbf F_p$ (21.5.2) et que $A$ est une algèbre sur $\mathbf F_p$ (de sorte
que $\Omega^1_A=\Omega^1_{A/\mathbf F_p}$), on déduit de (20.5.14) que l'on a
une \emph{suite exacte}
\begin{equation}
\label{0.22.4.6.1-fr}
  0\longrightarrow V\longrightarrow\Omega^1_A\otimes_AK
  \longrightarrow\Omega^1_K\longrightarrow0.
  \tag{22.4.6.1}
\end{equation}
On a donc par (22.4.5.3) un homomorphisme noté
\begin{equation}
\label{0.22.4.6.2-fr}
  \chi'_{B/A}:\Xi_{B_{(K)}/K}\longrightarrow V
  \tag{22.4.6.2}
\end{equation}
que nous appellerons \emph{homomorphisme caractéristique} de la $A$-algèbre
$B$ (relative à $i$, $j$ et $\mathfrak m$). Il résulte des définitions que pour
que $\chi'_{B/A}$ soit \emph{injectif}, il faut et il suffit que
$\pi'_{B_{(K)}/K}$ soit injectif, le noyau de ce dernier étant évidemment
contenu dans $\Xi_{B_{(K)}/K}$.
\end{env}

\begin{proposition}[22.4.7]
\label{0.22.4.7-fr}
Soient $A$ un anneau local artinien dont l'idéal maximal $\mathfrak m$ est de
carré nul, $K=A/\mathfrak m$ son corps résiduel, $V$ l'idéal $\mathfrak m$
considéré comme $K$-espace vectoriel~; on suppose que $A$ est de
caractéristique $p>0$~; soient
\begin{equation}
\label{0.22.4.7.1-fr}
  F_i(T_i)=T_i^p-f_i\qquad\text{avec }f_i\in A\quad(1\leqslant i\leqslant r)
  \tag{22.4.7.1}
\end{equation}
et désignons par $B$ l'anneau quotient de $A[T_1,\ldots,T_r]$ par l'idéal
engendré par les $r$ polynômes $F_i$. Alors~:
\begin{enumerate}
  \item[\rm(i)] $B$ est un anneau local.
  \item[\rm(ii)] Si $\mathfrak m'$ est l'idéal maximal de $B$,
  $K'=B/\mathfrak m'$ son corps résiduel, $V'$ le $B$-module
  $\mathfrak m'/\mathfrak m'^2$ considéré comme $K'$-espace vectoriel, alors,
  pour que $\operatorname{rg}_K(V)=\operatorname{rg}_{K'}(V')$, il faut et il suffit que l'homomorphisme
  caractéristique de $K$-espaces vectoriels
  \begin{equation}
  \label{0.22.4.7.2-fr}
    \chi'_{B/A}:\Xi_{B_{(K)}/K}\longrightarrow V
    \tag{22.4.7.2}
  \end{equation}
  (cf. (22.4.6.2)) soit injectif.
\end{enumerate}
\end{proposition}

\begin{proof}
Désignons encore par $a_i$ $(1\leqslant i\leqslant r)$ la classe de $f_i$
mod.~$\mathfrak m$~; les $a_i$ ne sont plus nécessairement $p$-indépendants
sur $K^p$, mais on peut toujours supposer que la sous-famille

